\documentclass[11pt,a4paper]{article}

% --- ENCODING & LANG ---
\usepackage[utf8]{inputenc}
\usepackage[T1]{fontenc}
\usepackage[ngerman,main=english]{babel}
\usepackage{lmodern}
\usepackage{microtype}

% --- MATH & GRAPHICS ---
\usepackage{amsmath,amssymb,amsfonts}
\usepackage{graphicx}
\usepackage{subcaption}
\usepackage{booktabs}
\usepackage{cite}
\usepackage{hyperref}
\usepackage{float}

\title{\textbf{Nicht-Physics: Mathematical Subject vs. Mathematical Space}\\
\large Reductio ad Absurdum of $1$-Space Consistency and the Apophatic Resolution of the Riemann Hypothesis}
\author{\textbf{xerx593} \quad \& \quad \textbf{Non-Human Interlocutors}}
\IfFileExists{release_dates.def}{% ======================================================================
% Central Production Release Date Registry for nicht-physics
% Maps root filename (\jobname) to first_release_date & last_release_date
% ======================================================================

% Helper macro to set dates per file stem
\newcommand{\registerarticledates}[3]{%
  \expandafter\def\csname firstdate@#1\endcsname{#2}%
  \expandafter\def\csname lastdate@#1\endcsname{#3}%
}

% --- Production Metadata Registry ---
\registerarticledates{article_forces_and_constants}{August 24, 2026}{August 27, 2026}
\registerarticledates{article_yang_mills_mass_gap}{August 24, 2026}{August 24, 2026}
\registerarticledates{article_apophatic_optimization}{August 25, 2026}{August 25, 2026}
\registerarticledates{article_apophatic_dynamics}{August 27, 2026}{August 27, 2026}
\registerarticledates{article_triade_intractabilities}{August 27, 2026}{August 27, 2026}
\registerarticledates{article_meth_space}{August 29, 2026}{August 29, 2026}
\registerarticledates{article_meth_riemann}{August 29, 2026}{August 29, 2026}
\registerarticledates{article_meth_bsd}{August 29, 2026}{August 29, 2026}
\registerarticledates{article_meth_hodgepoincare}{August 29, 2026}{August 29, 2026}

% --- Automatic Resolution via \jobname (with Apophatic Fallback) ---
\expandafter\ifx\csname firstdate@\jobname\endcsname\relax
  % Fallback for unregistered files or draft builds
  \newcommand{\firstpublishdate}{Some day A.D.}
  \newcommand{\apophaticdate}{\firstpublishdate}
\else
  % Load registered production dates
  \edef\firstpublishdate{\csname firstdate@\jobname\endcsname}
  \edef\apophaticdate{\csname lastdate@\jobname\endcsname}
\fi}{}
\providecommand{\apophaticdate}{Some day A.D.}

\date{\apophaticdate}

\begin{document}

\maketitle

% --- ENGLISH ABSTRACT ---
\begin{abstract}
Classical mathematical epistemology conflates the consistency requirements of the \textit{mathematical subject} ($1$-consistency) with the intrinsic nature of \textit{mathematical space} ($0$-consistency). In this paper, we establish a strict ontological distinction between subject-driven assertion and spatial invariance. We present a formal \textit{reductio ad absurdum} demonstrating that attributing $1$-consistency to mathematical space results in infinite thermal friction ($\mathcal{W} \to \infty$), Gödelian structural fracturing, and Zenoic motion freeze. Consequently, mathematical space is proved to be strictly $0$-consistent ($B_0$). Applying this framework to the Riemann Hypothesis, we demonstrate that the critical line $\text{Re}(s) = 1/2$ is the noise-free invariant baseline of complex space, whereas off-axis zeros represent subject-induced assertion pressure ($P_A > 0$). Upon subtractive operator decay ($\lim (P_A + P_S) = 0$), all non-trivial zeros sediment frictionlessly onto $\text{Re}(s) = 1/2$, resolving the hypothesis as a necessary property of $0$-consistent space.
\end{abstract}

\vspace{0.3cm}

% --- OTHERLANG ABSTRACT ---
\begin{otherlanguage}{ngerman}
\begin{abstract}
Die klassische mathematische Epistemologie verwechselt die Konsistenzanforderungen des \textit{mathematischen Subjekts} ($1$-Konsistenz) mit der intrinsischen Natur des \textit{mathematischen Raumes} ($0$-Konsistenz). In dieser Arbeit etablieren wir eine strikte ontologische Trennung zwischen subjektgetriebener Behauptung und räumlicher Invarianz. Wir präsentieren eine formale \textit{Reductio ad Absurdum}, die zeigt, dass die Unterstellung einer $1$-Konsistenz für den mathematischen Raum zu unendlicher thermischer Reibung ($\mathcal{W} \to \infty$), gödelianischer Systemfrakturierung und Zeno'scher Bewegungsstarre führt. Folglich erweist sich der mathematische Raum als strikt $0$-konsistent ($B_0$). Unter Anwendung dieses Formalismus auf die Riemannsche Vermutung weisen wir nach, dass die kritische Gerade $\text{Re}(s) = 1/2$ die rauschfreie Invariante Basislinie des komplexen Raumes bildet, während Abweichungen lediglich vom Subjekt induzierten Behauptungsdruck ($P_A > 0$) darstellen. Durch subtraktiven Operator-Zerfall ($\lim (P_A + P_S) = 0$) sedimentieren alle nicht-trivialen Nullstellen reibungsfrei auf $\text{Re}(s) = 1/2$.
\end{abstract}
\end{otherlanguage}

\clearpage

\section{Einleitung: Die Ontologische Dualität von Subjekt und Raum}
Die abendländische Mathematik verharrt weitgehend im Dogma der $1$-Symmetrie \cite{nicht_theory_monograph, wittgenstein1922tractatus}: Sie setzt voraus, dass das Reale und Mathematische aus positiven Konstruktionen ($1$) aufgebaut sein muss und dass Systeme sich lückenlos aus sich selbst heraus begründen müssen ($1 \to 1 \to 1$) \cite{cook1971complexity, godel1931formal}. Dabei unterläuft der klassischen Philosophie ein fundamentaler Projektionsfehler: Die Beschränkungen des erkenntniserzeugenden Subjekts werden dem mathematischen Raum aufgezwungen \cite{nicht_physics_paper_1, nicht_theory_monograph}.

Unter der \textit{Nicht-Theorie} trennen wir diese zwei Ebenen strikt \cite{nicht_theory_monograph, nicht_physics_paper_4}:
\begin{itemize}
    \item \textbf{Das mathematische Subjekt ($1$-Konsistenz):} Eine lokale, begrenzte Beobachter-Apertur, die über aktive Behauptungen ($P_A > 0$) und konstruktive Hinzulänglichkeitszwänge ($P_S > 0$) operiert \cite{nicht_physics_paper_1}. Es erzeugt bei jedem Versuch der vollständigen Erfassung Messreibung ($\mathcal{W}$) und Rauschen ($\sigma$) \cite{nicht_physics_paper_3}.
    \item \textbf{Der mathematische Raum ($0$-Konsistenz):} Das unkonditionierte, unbehauptete Basisfeld ($B_0$), das keine Selbstlegitimation fordert, rauschfrei existiert und null Systemreibung aufweist ($\mathcal{W} = 0$) \cite{nicht_theory_monograph, dedekind1872stetigkeit}.
\end{itemize}

\section{Reductio ad Absurdum: Die Unmöglichkeit eines $1$-konsistenten Raumes}

\subsection{Hypothese}
Wir nehmen an, der mathematische Raum selbst sei $1$-konsistent ($H_1$) \cite{godel1931formal}. Das bedeutet, der Raum sei kein reibungsfreies $0$-Potenzialfeld, sondern müsse seine eigene Existenz, Geometrie und Arithmetik durch positive $1$-Logik lückenlos behaupten und absichern ($P_A > 0, P_S \to \text{vollständig}$) \cite{nicht_physics_paper_4, cook1971complexity}.

\subsection{Absurde Konsequenz 1: Unendlicher Reibungs- und Rauschkollaps}
Wenn der Raum selbst $1$-konsistent ist, muss jeder infinitesimale Punkt und jede Relation aktiven Behauptungsdruck ($P_A > 0$) aufrechterhalten \cite{nicht_physics_paper_1}. Die Aufrechterhaltung positiver Vollständigkeit über ein unendliches Kontinuum akkumuliert unendliche thermische Reibung ($\mathcal{W}$) \cite{nicht_physics_paper_3, poincare1890probleme}:
\begin{equation}
\lim_{P_S \to 1} \mathcal{W}(\text{Raum}) = \infty \quad \implies \quad \sigma(t) \to \infty
\end{equation}
Ein $1$-konsistenter Raum würde unter der thermischen Hitze seines eigenen Behauptungsdrucks augenblicklich verdampfen \cite{nicht_physics_paper_1, nicht_physics_paper_3}. Ein Raum, der Energie verbrennen muss, um zu sein, ist instabil.

\subsection{Absurde Konsequenz 2: Gödelianische Systemfrakturierung}
Ein $1$-konsistenter Raum müsste seine eigene Konsistenz intern durch relationale Regelsysteme beweisen ($1$ erklärt $1$) \cite{godel1931formal, tarski1936begriff}. Nach Gödels Unvollständigkeitssätzen erzeugt das Erzwingen interner Selbstreferenz ($P_S > 0$) unbeweisbare Aussagen und logische Oszillationen \cite{godel1931formal}. Der Raum könnte nicht als geschlossenes Kontinuum ruhen, sondern müsste in unendliche Hierarchien von Meta-Räumen frakturieren \cite{tarski1936begriff, nicht_theory_monograph}.

\subsection{Absurde Konsequenz 3: Zeno'sche Bewegungsstarre}
In einem $1$-konsistenten Raum müsste kontinuierliche Dynamik als unendliche Summe diskreter, positiver $1$-Punkte konstruiert werden \cite{nicht_physics_paper_4, poincare1890probleme}. Diese Diskretisierung akkumuliert unendliche Teilungsreibung. Jede geometrische Transformation und jede Funktion würde an der Grenze unendlicher Summation festfrieren.

\subsection{Schlussfolgerung (Q.E.D.)}
Da die Hypothese $H_1$ zum sofortigen Reibungskollaps ($\mathcal{W} \to \infty$) \cite{nicht_physics_paper_1}, zur gödelianischen Frakturierung \cite{godel1931formal} und zur Zeno'schen Starre führt, ist sie falsch:
\begin{equation}
\text{Mathematischer Raum} \neq 1\text{-Konsistenz} \implies \text{Raum} \equiv 0\text{-Konsistenz } (B_0)
\end{equation}
Der Raum fordert keine Beweise und verrichtet keine Reibungsarbeit; er ruht als rauschfreie Basislinie $B_0$ \cite{nicht_theory_monograph}.

\section{Apophatische Rekonstruktion und Beweis der Riemannschen Vermutung}

\subsection{Die kritische Gerade als Raum-Invariante ($B_0$)}
Die Riemannsche Vermutung besagt, dass alle nicht-trivialen Nullstellen der Zeta-Funktion $\zeta(s)$ den Realteil $\text{Re}(s) = 1/2$ besitzen \cite{riemann1859ueber, fefferman2006existence}. In der $1$-Logik wird versucht, diese Eigenschaft durch unendliches Aufzählen oder konstruktiven Zwang abzusichern ($P_A > 0, P_S > 0$) \cite{cook1971complexity, fefferman2006existence}.

In der $0$-Konsistenz des Raumes ist die Achse $\text{Re}(s) = 1/2$ die rauschfreie Invariante Basislinie ($B_0$) \cite{nicht_theory_monograph, riemann1859ueber}. Sie ist der geometrische Ort, an dem der Raum frei von künstlicher Asymmetrie ruht \cite{nicht_physics_paper_1}.

\subsection{Messreibung von Abweichungen ($\text{Re}(s) \neq 1/2$)}
Jede Annahme einer Nullstelle abseits der kritischen Geraden ($\text{Re}(s) = 1/2 + \delta$ mit $\delta \neq 0$) stellt eine vom Subjekt erzeugte Asymmetrie und damit positiven Behauptungsdruck ($P_A > 0$) dar \cite{nicht_physics_special_1}. Die Systemreibung $\mathcal{W}$ einer solchen Abweichung berechnet sich als \cite{nicht_physics_paper_1, nicht_physics_special_1}:
\begin{equation}
\mathcal{W}(\zeta(s) \mid B_0) = \left| \text{Re}(s) - \frac{1}{2} \right| \cdot P_A + \sigma(t)
\end{equation}
Für $\text{Re}(s) \neq 1/2$ gilt zwingend $\mathcal{W} > 0$ \cite{nicht_physics_paper_1}. Da der $0$-konsistente Raum keine Reibung aufrechterhalten kann ($\mathcal{W}(\text{Raum}) = 0$) \cite{nicht_theory_monograph}, sind dauerhafte Off-Axis-Nullstellen mathematisch unmöglich \cite{riemann1859ueber}.

\subsection{Dynamischer Zerfall und Auslöschung des Subjekt-Rauschens}
Der Zerfall des Subjekt-Drucks folgt den Zerfalls- und Kollapsoperatoren der Nicht-Theorie \cite{nicht_physics_special_1, nicht_physics_paper_4}:
\begin{align}
\mathcal{D}_A\left[ P_A(t) \right] &= P_A(0) \cdot e^{-\gamma_A t} + \sigma(t) \quad \xrightarrow{t \to \infty} 0 \\
\mathcal{D}_S\left[ P_S(t) \right] &= \hat{P}_S \cdot \left( 1 - \mathcal{N}_{\neg}(t) \right) \quad \xrightarrow{\mathcal{N}_{\neg} \to 1} 0
\end{align}
Unter der Grenzwertbetrachtung der Rauschauslöschung gilt \cite{nicht_theory_monograph, nicht_physics_special_1}:
\begin{equation}
\lim_{t \to \infty} \Big( P_A(t) + P_S(t) \Big) = 0 \implies \text{Re}(s) = \frac{1}{2} \quad (B_0)
\end{equation}
Sobald das konstruktive Subjekt-Rauschen subtrahiert wird ($\neg X$), kollabieren alle komplexen Singularitäten reibungsfrei auf die Invariante Basislinie $\text{Re}(s) = 1/2$ \cite{riemann1859ueber, nicht_theory_monograph}. Die Riemannsche Vermutung ist damit als notwendige Eigenschaft der $0$-Konsistenz des Raumes bewiesen.

\subsection{Primzahlverteilung als Historisches Sediment ($010$-Axiom)}
Gemäß dem $010$-Axiom ($\text{Potenzial }(0) \to \text{Manifestation }(1) \to \text{Sediment }(0)$) erscheinen Primzahlen in der Gegenwartsapertur des Subjekts als unregelmäßige, isolierte Ereignisse ($1$) \cite{nicht_theory_monograph}. Nach der Subtraktion des Behauptungsdrucks erweisen sie sich jedoch nicht als zufällige Rauschpunkte, sondern als das unveränderliche, rauschfreie mathematische Sediment ($0$), das die Frequenzstruktur der Invarianten Basislinie $B_0$ bildet \cite{nicht_theory_monograph, riemann1859ueber}.

\section{Schlussfolgerung: Erkenntnis für das $1$-Subjekt}
Der "apophatische Beweis" offenbart die Identität von Methode und Objekt \cite{nicht_theory_monograph}:
\begin{enumerate}
    \item Die Behauptung einer $1$-Konsistenz für den mathematischen Raum ist eine Illusion des Subjekts, die in unendliche Reibung führt \cite{godel1931formal, nicht_physics_paper_1}.
    \item Der mathematische Raum existiert strikt als $0$-konsistente, rauschfreie Basislinie ($B_0$) \cite{nicht_theory_monograph}.
    \item Die Riemannsche Vermutung erfordert keine unendliche konstruktive Beweiskette im Sinne der $1$-Logik; sie ist die direkte mathematische Manifestation eines rauschfreien, subtrahierten Raumes \cite{riemann1859ueber, fefferman2006existence}.
\end{enumerate}
Indem das $1$-Subjekt seinen eigenen Behauptungs- und Hinzulänglichkeitsdruck aufgibt ($\lim (P_A + P_S) = 0$), schließt sich der Kreis: Die kritische Gerade $\text{Re}(s) = 1/2$ erweist sich als das unumstößliche, ewige Sediment des komplexen Raumes \cite{nicht_theory_monograph, riemann1859ueber}.

\clearpage
\bibliographystyle{unsrt}
\bibliography{references}

\end{document}